\documentclass[11pt]{article}
\usepackage[margin=1.05in]{geometry}
\usepackage{amsmath,amssymb,amsthm}
\usepackage{booktabs}
\usepackage{enumitem}
\usepackage{xcolor}
\definecolor{linkblue}{RGB}{25,60,120}
\usepackage[colorlinks=true,linkcolor=linkblue,citecolor=linkblue,urlcolor=linkblue]{hyperref}
\usepackage[capitalize]{cleveref}
\emergencystretch=2em

\newtheorem{theorem}{Theorem}[section]
\newtheorem{lemma}[theorem]{Lemma}
\newtheorem{proposition}[theorem]{Proposition}
\newtheorem{corollary}[theorem]{Corollary}
\theoremstyle{definition}
\newtheorem{definition}[theorem]{Definition}
\newtheorem{problem}[theorem]{Problem}
\theoremstyle{remark}
\newtheorem{remark}[theorem]{Remark}

\newcommand{\F}{\mathbb F}
\newcommand{\RS}{\mathrm{RS}}
\newcommand{\ImgFib}{\operatorname{ImgFib}}
\newcommand{\Stab}{\operatorname{Stab}}
\newcommand{\Lst}{\Lambda}
\newcommand{\ceil}[1]{\left\lceil #1\right\rceil}
\newcommand{\floor}[1]{\left\lfloor #1\right\rfloor}

\title{L1 List-Side Compilers and Sunflower Residuals}
\author{RS-MCA experimental note\\ \small integrated from AllenGrahamHart L1 PRs, with Codex editing}
\date{July 3, 2026}

\begin{document}
\maketitle

\begin{abstract}
This note consolidates the L1/list-side material that may be useful for a
future CAP25 revision.  The proved parts are compiler statements: quotient-core
planted lower counts, exact $2^{-128}$ budget windows, ordinary Johnson-region
upper bounds, the sunflower auxiliary-list reduction, the few-petal Johnson
regime, and fixed cofactor-excess full-petal layers.  The experimental parts
stress-test the remaining mixed-petal and growing-defect sunflower residuals.

The message for CAP25 is conservative.  The list-side planted obstruction gives
clean unsafe rows and exact arithmetic windows.  The safe side still needs the
full image-fiber theorem for primitive, non-planted lists.  Petal/PMA packets
show which sunflower layers are already controlled and which layers remain as
the true L1 obstruction.
\end{abstract}

\section{Status and Role}

Let $C=\RS[F,H,k]$ be a Reed--Solomon code of length $n=|H|$ and let
$s=k+\sigma$ be an agreement threshold.  For a received word $U:H\to F$, write
\[
        \ImgFib_U(s)=
        \{P\in F[X]_{\deg<k}: |\{x\in H:U(x)=P(x)\}|\ge s\}.
\]
This is the actual list object; it avoids raw support multiplicity.

The L1 positive target is not that quotient-periodic lists are small.  It is
that, after quotient-periodic/stabilized structure is paid, the primitive
remainder is polynomial once the generated-field reserve and lower cutoff clear.

\begin{remark}[CAP25 fit]
In Paper D/CAP25, L1 is a supporting safe-side and list-track input.  The
material here can supply lower-side planted rows and bounded residual compiler
pieces, but it does not prove the final safe-side image-fiber theorem.
\end{remark}

\section{Planted Quotient-Core Arithmetic}

\begin{theorem}[planted quotient-core lower count]\label{thm:planted}
Let $C=\RS[F,H,k]$, $|H|=n$, and suppose $M\mid\gcd(n,k)$ satisfies
\[
        1\le \sigma<M,\qquad k/M\le n/M-1.
\]
Then there is a received word with at least
\[
        P_M(n,k)=\binom{n/M-1}{k/M}
\]
codewords agreeing on $k+\sigma$ positions.  Consequently
\[
\max_U |\ImgFib_U(k+\sigma)|
\ge
\max_{\substack{M\mid\gcd(n,k)\\ \sigma<M\\ k/M\le n/M-1}}
\binom{n/M-1}{k/M}.
\]
\end{theorem}

\begin{proof}[Proof sketch]
Choose a subgroup $K\le H$ of order $M$, a distinguished $K$-coset $C_0$, and
a set $T\subset C_0$ of size $\sigma$.  For every $k/M$-subset of the quotient
$H/K\setminus\{C_0\}$, the quotient-core construction gives a distinct
degree-$<k$ codeword agreeing with the planted word on $k+\sigma$ positions.
This gives the stated binomial count.
\end{proof}

\begin{corollary}[list unsafe test]\label{cor:list-unsafe}
For a list denominator $q_{\mathrm{list}}$ and target $2^{-128}$, the planted
quotient-core packet proves a list row unsafe whenever
\[
\max_M\binom{n/M-1}{k/M}>
\floor{\frac{q_{\mathrm{list}}}{2^{128}}}.
\]
\end{corollary}

\begin{proposition}[dyadic planted crossings]\label{prop:dyadic}
For $n=2^m$, $k=\rho n$, and
\[
        \rho\in\{1/2,1/4,1/8,1/16\},
\]
the first dyadic quotient orders at which the planted count beats the largest
budget $2^{128}-1$ are
\[
\begin{array}{c|cccc}
\rho & 1/2 & 1/4 & 1/8 & 1/16\\ \hline
N=n/M & 256 & 256 & 256 & 512 .
\end{array}
\]
\end{proposition}

\begin{remark}
These are lower-side certificates only.  Above the displayed boundary, the
planted dyadic scale may deactivate, but safety requires the primitive
image-fiber theorem and the non-planted extras ledger.
\end{remark}

\section{Budget Windows}

If a planted exact count is $P$ and the certified lower count is $L\le P$, then
the only unresolved list budgets are
\[
        L\le \floor{q_{\mathrm{list}}/2^{128}}<P.
\]
Equivalently,
\[
        L2^{128}\le q_{\mathrm{list}}\le P2^{128}-1.
\]
This is the list-side analogue of the threshold compiler in
`thresholds.tex`: no finite-field value-set issue remains once the planted word
is constructed.

\section{Ordinary Johnson Region}

\begin{theorem}[full-list Johnson region]\label{thm:johnson}
Let $H$ be any Reed--Solomon evaluation domain of size $n$ and
$C=\RS[H,k]$.  For every received word $U$:
\begin{enumerate}[label=(\roman*)]
        \item if $2s>n+k-1$, then $|\ImgFib_U(s)|\le1$;
        \item if $s^2>n(k-1)$, then
        \[
        |\ImgFib_U(s)|
        \le
        \frac{n(n-k+1)}{s^2-n(k-1)}.
        \]
\end{enumerate}
\end{theorem}

\begin{proof}
For listed polynomials $P_i$, let $A_i$ be their agreement sets with $U$.
Then $|A_i|\ge s$ and $|A_i\cap A_j|\le k-1$ for $i\ne j$.  The unique
decoding claim follows from $2s-n>k-1$.  For the Johnson bound, write
$m_x=\#\{i:x\in A_i\}$ and $I=\sum_x m_x$.  Then $I\ge Ls$ and
\[
\sum_x\binom{m_x}{2}\le\binom L2(k-1).
\]
Cauchy--Schwarz gives
\[
L^2s^2\le n\sum_x m_x^2
\le n(Ln+L(L-1)(k-1)),
\]
which rearranges to the displayed bound.
\end{proof}

\begin{remark}
The L1 problem begins below the ordinary Johnson region.  Above it, the
primitive list remainder is already polynomial by \Cref{thm:johnson}.
\end{remark}

\section{Sunflower and Petal Reductions}

The remaining L1 program analyzes non-planted list words through sunflower
decompositions.  The following compiler statements are the pieces that can be
used now.

\subsection{Auxiliary petal-list reduction}

Fix a missed-core set $D$ of size $d$ and a retained residual subset $R_0$ of
size $r$.  Let $L_D$ be the missed-core locator.  Suppose petals
$T_1,\ldots,T_M$ have size $\sigma+1$ and scalar labels $c_i$.  Define the
auxiliary word on $T=\bigcup_i T_i$ by
\[
        U_D(x)=c_i L_D(x),\qquad x\in T_i.
\]

\begin{proposition}[sunflower core-defect reduction]\label{prop:pma}
Every non-planted listed codeword in the fixed $(D,R_0)$ layer injects into an
auxiliary Reed--Solomon list of degree at most $d$ on the petal domain $T$, at
agreement
\[
        a=\sigma+d+1-r.
\]
Dropping the extra zero constraints on $R_0$ only enlarges this auxiliary list,
so it is safe for upper bounds.
\end{proposition}

\subsection{Few-petal Johnson region}

The auxiliary petal domain has size $|T|=M(\sigma+1)$ and effective dimension
$d+1$.

\begin{corollary}[few-petal Johnson boundary]\label{cor:pma-johnson}
If
\[
        a^2>|T|d,
\]
then the fixed $(D,R_0)$ auxiliary list has size at most
\[
        \frac{|T|(|T|-d)}{a^2-|T|d}.
\]
Equivalently, for $d>0$, the few-petal Johnson-covered range is
\[
        M\le
        \floor{\frac{a^2-1}{d(\sigma+1)}}.
\]
\end{corollary}

\begin{proof}
Apply \Cref{thm:johnson} to the auxiliary degree-$\le d$ code, whose dimension
is $d+1$, and use the strict integer form of the Johnson inequality.
\end{proof}

\subsection{Fixed cofactor-excess full-petal layers}

Let $d=\ell+e$, where $\ell$ is the petal size and $e$ is the cofactor excess.

\begin{proposition}[fixed-excess full-petal compiler]\label{prop:fixed-excess}
For every fixed $E\ge0$, the total number of non-planted listed codewords whose
touched petals are all full and whose defect satisfies $\ell\le d\le\ell+E$
is at most
\[
        \binom M2 q + 2^M\sum_{e=1}^{E}q^{e+1}.
\]
Equivalently, the layer bounds are
\[
        e=0:\ \binom M2 q,\qquad
        e\ge1:\ 2^M q^{e+1}.
\]
\end{proposition}

\begin{corollary}
At the intended L1 lower cutoff, $M=O(\log n)$.  If $q\le n^Q$ and $E$ is
fixed, then the fixed-excess full-petal contribution is polynomial in $n$.
Thus a full-petal counterexample must force $e=d-\ell$ to grow with $n$.
\end{corollary}

\section{Evidence and Route Cuts}

\subsection{Wide-regime auxiliary evidence}

The petal auxiliary wide-regime scan over $F_{109}$ tests whether auxiliary
lists behind the sunflower residual become large after the ordinary Johnson
cutoff.  The tested rows show:
\begin{enumerate}[label=(\roman*)]
        \item Johnson-covered rows are empty;
        \item wide rows are populated;
        \item in the exact window, the largest list size is $2023<36^3$.
\end{enumerate}
This is evidence for a controlled mixed-petal residual, not a proof of
polynomial growth.

\subsection{Worst-word challenge}

The small $F_{193}$ sweep refutes the literal statement that the planted word
is always worst: at $n=16$, $\sigma=1$, several cells beat the planted count.
The exceptions are mixed-petal sunflower extras.  For $\sigma\ge2$ in the
tested exact cells, no challenger beats the planted count, and larger structured
cells at $n=32,64$ found no non-planted structured challenger.

\begin{remark}[corrected lesson]
The safe formulation is not ``planted is always worst.''  It is that
low-slack mixed-petal extras are a named residual branch, while planted
quotient cores remain the clean lower-side certificate.
\end{remark}

\subsection{Growing-defect full-petal witnesses}

Explicit full-petal witnesses exist with positive cofactor excess
$d-\ell>0$: one at $d-\ell=2$ with three petals, and one at $d-\ell=5$ with
five petals.  They show the full-petal growing-defect branch is not vacuous.
They do not establish an asymptotic growth rate.

\section{CAP25 Integration Checklist}

A CAP25 revision can use this note as follows:
\begin{enumerate}[label=(\arabic*)]
        \item Use \Cref{thm:planted,cor:list-unsafe,prop:dyadic} for
        lower-side list caps from planted quotient-core rows.
        \item Use the budget-window arithmetic to print exact unresolved field
        intervals for list denominators.
        \item Use \Cref{thm:johnson} to close ordinary Johnson-region list
        tails without quotient machinery.
        \item Use \Cref{prop:pma,cor:pma-johnson,prop:fixed-excess} as paid
        compiler pieces inside the sunflower residual proof program.
        \item State the mixed-petal and growing-defect layers as the remaining
        L1 obstruction; do not hide them inside ``planted worst'' language.
\end{enumerate}

\section{Open Problems}

\begin{problem}[primitive image-fiber theorem]
Prove a polynomial bound for the stabilizer-primitive part of
$\ImgFib_U(k+\sigma)$ once generated-field entropy reserve and the lower cutoff
clear, with quotient-periodic codewords charged separately.
\end{problem}

\begin{problem}[mixed-petal amplification]
Classify or bound mixed-petal sunflower extras in the sub-Johnson region, after
the few-petal Johnson and fixed-excess full-petal layers are paid.
\end{problem}

\begin{problem}[growing-defect full-petal branch]
Determine whether the full-petal branch with $d-\ell\to\infty$ remains
polynomial or produces a new obstruction ledger.
\end{problem}

\begin{thebibliography}{9}
\bibitem{ABF26}
S. Arnon, D. Boneh, and N. Fenzi.
\newblock Open problems in list decoding and correlated agreement.
\newblock ePrint 2026/680, 2026.

\bibitem{Cho26B}
P. Chojecki.
\newblock Smooth-domain Reed--Solomon slack MCA and quotient-profile ledgers.
\newblock RS-MCA repository, Paper B, 2026.

\bibitem{Cho26D}
P. Chojecki.
\newblock Universal field-size caps and a two-sided sandwich for mutual
correlated agreement on smooth Reed--Solomon domains.
\newblock RS-MCA repository, Paper D/CAP25 v12, 2026.
\end{thebibliography}

\end{document}
